\chapter{Dynamic Kernel Extensions}
\label{chap:kernel}

The Linux kernel is a monolith, but it is a \emph{programmable} monolith. AULinux
deliberately ships an unmodified upstream kernel and layers its distinctive
behaviour on top through three small, self-contained Loadable Kernel Modules
(LKMs): \file{aulinux\_core}, \file{aulinux\_procmon}, and \file{aulinux\_sysctl}.
Each is an out-of-tree module compiled against the running kernel's headers and
loaded at boot. This chapter dissects what those modules \emph{actually} do —
line by line — and is explicit about which parts are production-ready and which
are scaffolding awaiting completion.

\begin{infobox}[The three modules at a glance]
\code{aulinux\_core} surfaces version telemetry under \file{/proc/aulinux};
\code{aulinux\_procmon} counts process lifecycle events via kprobes and exposes
them under \file{/proc/aulinux/procmon}; \code{aulinux\_sysctl} registers runtime
tunables under \file{/proc/sys/aulinux/} and exports two symbols that the other
modules consume. All three are GPL-licensed and authored by ``AULinux Team''.
\end{infobox}

\section{Why Modules Instead of Patching the Kernel}

There are two ways to add behaviour to Linux: patch the source tree and ship a
custom kernel, or write loadable modules that bolt onto a stock kernel at
runtime. AULinux chooses the latter, and the rationale is worth stating plainly.

\begin{itemize}
  \item \textbf{Decoupling from the upstream release cadence.} A forked kernel
        must be re-patched and re-tested against every upstream security fix. An
        out-of-tree module only needs to compile against the kernel's exported
        headers, so AULinux can ride stock kernel updates with far less churn.
  \item \textbf{Smaller blast radius.} A bug in a module that fails to load
        leaves a bootable system; a bug baked into \code{vmlinux} may not.
  \item \textbf{Runtime composability.} Modules can be \code{insmod}'d and
        \code{rmmod}'d independently, which is exactly how the AULinux build and
        test scripts exercise them.
\end{itemize}

\begin{notebox}[The module ABI is only ``stable-ish'']
Linux offers \emph{no} stable in-kernel ABI. Symbols, struct layouts, and even
the signatures of hookable functions change between releases. Out-of-tree
modules pay for their decoupling by being recompiled against each new kernel —
and, as we will see with \code{aulinux\_procmon}, by occasionally targeting
kernel internals (like \code{do\_fork}) that have been renamed or removed
upstream.
\end{notebox}

\section{The Shared Interface: \file{aulinux.h}}

All three modules are meant to share a single header, \file{aulinux.h}, which
defines the project's common vocabulary. It pulls in only \code{<linux/types.h>}
and is guarded by the conventional \code{\_AULINUX\_H} include guard.

\begin{lstlisting}[language=C, caption={Versioning and feature flags from aulinux.h}]
/* Version information */
#define AULINUX_VERSION_MAJOR  1
#define AULINUX_VERSION_MINOR  0
#define AULINUX_VERSION_PATCH  0
#define AULINUX_VERSION_STRING "1.0.0"

/* Feature flags (for aulinux_feature_enabled) */
#define AULINUX_FEATURE_PROCMON     (1 << 0)  /* Process monitoring */
#define AULINUX_FEATURE_AUDIT       (1 << 1)  /* Security auditing */
#define AULINUX_FEATURE_SANDBOX     (1 << 2)  /* Process sandboxing */
#define AULINUX_FEATURE_NETMON      (1 << 3)  /* Network monitoring */
\end{lstlisting}

The header's most load-bearing contribution is the cross-module contract. It
declares two functions as \code{extern} and provides logging macros that gate
debug output on a runtime level:

\begin{lstlisting}[language=C, caption={Cross-module declarations and level-gated logging (aulinux.h, lines 30--48)}]
/* External function declarations (from aulinux_sysctl.c) */
extern int aulinux_get_debug_level(void);
extern int aulinux_feature_enabled(int flag);

/* Logging macros with debug level checking */
#define aulinux_pr_debug(level, fmt, ...) \
    do { \
        if (aulinux_get_debug_level() >= (level)) \
            pr_debug("AULinux: " fmt, ##__VA_ARGS__); \
    } while (0)

#define aulinux_pr_info(fmt, ...)  pr_info("AULinux: " fmt, ##__VA_ARGS__)
#define aulinux_pr_warn(fmt, ...)  pr_warn("AULinux: " fmt, ##__VA_ARGS__)
#define aulinux_pr_err(fmt, ...)   pr_err("AULinux: " fmt, ##__VA_ARGS__)
\end{lstlisting}

The header also reserves an IOCTL command space (magic byte \code{'A'} with
\code{GET\_VERSION}, \code{GET\_FEATURES}, \code{SET\_FEATURES}), a set of size
limits (\code{AULINUX\_MAX\_NAME\_LEN}, \code{AULINUX\_MAX\_PATH\_LEN},
\code{AULINUX\_MAX\_PROCESSES}), and a small return-code enumeration
(\code{AULINUX\_SUCCESS} through \code{AULINUX\_ERR\_BUSY}).

\begin{notebox}[Header is aspirational in part]
The IOCTL command numbers, the feature-flag bitfield names, and the
\code{AULINUX\_ERR\_*} return codes are \emph{declared but not yet used} by any
of the three modules. No module currently registers a character device or
implements an \code{.unlocked\_ioctl} handler, so the IOCTL interface is a
forward-looking placeholder. Only \code{aulinux\_get\_debug\_level()} and
\code{aulinux\_feature\_enabled()} are wired up — and notably, \file{aulinux.h}
is \code{\#include}d only by \file{aulinux\_sysctl.c}; the other two modules do
not yet consume the shared logging macros.
\end{notebox}

\section{\code{aulinux\_core}: The Telemetry Anchor}

\file{aulinux\_core.c} is the foundational module. Its single responsibility is
to establish the \file{/proc/aulinux} directory and publish a read-only
\file{version} file inside it. It is small, correct, and complete.

\subsection{Module Metadata}

The module identifies itself with the standard \code{MODULE\_*} macros. The
license string \code{"GPL"} is significant: it grants the module access to
GPL-only exported kernel symbols.

\begin{lstlisting}[language=C, caption={Module metadata (aulinux\_core.c, lines 16--22)}]
#define AULINUX_VERSION "1.0.0"
#define AULINUX_PROC_NAME "aulinux"

MODULE_LICENSE("GPL");
MODULE_AUTHOR("AULinux Team");
MODULE_DESCRIPTION("AULinux Core Kernel Module");
MODULE_VERSION(AULINUX_VERSION);
\end{lstlisting}

\subsection{The \code{proc\_ops} / \code{seq\_file} Pattern}

Modern kernels (5.6+) replaced the old \code{file\_operations} struct for
\file{/proc} entries with a leaner \code{struct proc\_ops}. AULinux uses it
together with the \code{seq\_file} helper API, which is the idiomatic way to emit
small, formatted, read-only text from the kernel without manually juggling user
buffers and offsets.

\begin{lstlisting}[language=C, caption={The version producer and its proc\_ops (aulinux\_core.c, lines 31--51)}]
static int aulinux_version_show(struct seq_file *m, void *v)
{
    seq_printf(m, "AULinux Version: %s\n", AULINUX_VERSION);
    seq_printf(m, "Kernel Version: %d.%d.%d\n",
               LINUX_VERSION_MAJOR,
               LINUX_VERSION_PATCHLEVEL,
               LINUX_VERSION_SUBLEVEL);
    return 0;
}

static int aulinux_version_open(struct inode *inode, struct file *file)
{
    return single_open(file, aulinux_version_show, NULL);
}

static const struct proc_ops aulinux_version_ops = {
    .proc_open    = aulinux_version_open,
    .proc_read    = seq_read,
    .proc_lseek   = seq_lseek,
    .proc_release = single_release,
};
\end{lstlisting}

The flow is worth tracing because every \file{/proc} entry in the AULinux
modules follows it. When userspace \code{open()}s the file, the kernel calls
\code{aulinux\_version\_open}, which registers \code{aulinux\_version\_show} as a
single-shot producer via \code{single\_open()}. On \code{read()}, the generic
\code{seq\_read} invokes the show function, which streams formatted lines into a
managed buffer; \code{single\_release} tears it down. The reported kernel version
comes from compile-time macros in \code{<linux/version.h>}, so it reflects the
headers the module was \emph{built} against.

\subsection{Init and Exit Paths}

The lifecycle of every LKM is bracketed by a function tagged \code{\_\_init}
(reclaimed from memory after load) and one tagged \code{\_\_exit}, registered via
\code{module\_init()} / \code{module\_exit()}. The core module's init path builds
its \file{/proc} tree and unwinds carefully on failure:

\begin{lstlisting}[language=C, caption={Init path with ordered cleanup (aulinux\_core.c, lines 56--79)}]
static int __init aulinux_core_init(void)
{
    pr_info("AULinux: Initializing core module v%s\n", AULINUX_VERSION);

    aulinux_proc_dir = proc_mkdir(AULINUX_PROC_NAME, NULL);
    if (!aulinux_proc_dir) {
        pr_err("AULinux: Failed to create /proc/%s\n", AULINUX_PROC_NAME);
        return -ENOMEM;
    }

    aulinux_version_entry = proc_create("version", 0444,
                                         aulinux_proc_dir,
                                         &aulinux_version_ops);
    if (!aulinux_version_entry) {
        pr_err("AULinux: Failed to create version proc entry\n");
        proc_remove(aulinux_proc_dir);   /* unwind on failure */
        return -ENOMEM;
    }

    pr_info("AULinux: Core module initialized successfully\n");
    return 0;
}
\end{lstlisting}

The exit path mirrors this, removing the \file{version} entry and then the
directory. Returning a negative errno from an \code{\_\_init} function aborts the
load, so this ordered cleanup is essential — leaving a half-built \file{/proc}
tree behind would leak kernel resources.

\begin{infobox}[What you actually see]
\begin{verbatim}
$ cat /proc/aulinux/version
AULinux Version: 1.0.0
Kernel Version: 6.8.0
\end{verbatim}
The mode \code{0444} makes the file world-readable but unwritable: pure
telemetry, no control surface.
\end{infobox}

\section{\code{aulinux\_procmon}: Watching Process Lifecycle}

\file{aulinux\_procmon.c} is the most ambitious module. Its goal is to observe
process creation and termination and expose both running counters and a recent-
event history, so that a userland \code{ps}-like tool could read process activity
straight from \file{/proc}.

\subsection{Mechanism: Kprobes on Kernel Internals}

The module uses \textbf{kprobes} — a kernel facility that lets you attach a
handler to (almost) any kernel symbol by patching a breakpoint at its entry. Two
probes are declared, each with a \code{pre\_handler} that fires \emph{before} the
target function runs:

\begin{lstlisting}[language=C, caption={Kprobe definitions targeting fork and exit (aulinux\_procmon.c, lines 104--112)}]
static struct kprobe fork_kp = {
    .symbol_name = "do_fork",
    .pre_handler = fork_handler_pre,
};

static struct kprobe exit_kp = {
    .symbol_name = "do_exit",
    .pre_handler = exit_handler_pre,
};
\end{lstlisting}

Each handler simply bumps an atomic counter and records an event against the
\code{current} task — the task descriptor of whatever process is executing the
probed code:

\begin{lstlisting}[language=C, caption={Pre-handlers increment counters and log events (aulinux\_procmon.c, lines 86--101)}]
static int fork_handler_pre(struct kprobe *p, struct pt_regs *regs)
{
    atomic_long_inc(&total_forks);
    record_event(PROCMON_FORK, current);
    return 0;
}

static int exit_handler_pre(struct kprobe *p, struct pt_regs *regs)
{
    atomic_long_inc(&total_exits);
    record_event(PROCMON_EXIT, current);
    return 0;
}
\end{lstlisting}

\begin{warnbox}[\code{do\_fork} no longer exists on modern kernels]
The fork probe targets the symbol \code{"do\_fork"}, which was renamed to
\code{\_do\_fork} and later removed in favour of \code{kernel\_clone()} around
Linux 5.10. On any recent kernel \code{register\_kprobe(\&fork\_kp)} will fail
with \code{-EINVAL} (no such symbol). The module is written to \emph{tolerate}
this: its init path logs a warning and continues, so fork counting will silently
read zero rather than crashing the system. To actually count forks on a current
kernel, the \code{.symbol\_name} must be updated to \code{"kernel\_clone"}.
\end{warnbox}

\subsection{The Event Ring Buffer}

Events are stored in a fixed 64-slot circular buffer protected by an
IRQ-safe spinlock — appropriate because probe handlers can fire in atomic
context. Each record captures the event type, PID, parent PID, the 16-byte
\code{comm} name, and a real-time timestamp.

\begin{lstlisting}[language=C, caption={Recording an event under an IRQ-safe lock (aulinux\_procmon.c, lines 58--81)}]
static void record_event(enum procmon_event_type type,
                         struct task_struct *task)
{
    unsigned long flags;

    if (!monitoring_enabled)
        return;

    spin_lock_irqsave(&event_lock, flags);

    event_history[event_head].type = type;
    event_history[event_head].pid  = task->pid;
    event_history[event_head].ppid = task->real_parent ?
                                     task->real_parent->pid : 0;
    strncpy(event_history[event_head].comm, task->comm, TASK_COMM_LEN - 1);
    event_history[event_head].comm[TASK_COMM_LEN - 1] = '\0';
    event_history[event_head].timestamp = ktime_get_real();

    event_head = (event_head + 1) % PROCMON_HISTORY_SIZE;
    if (event_count < PROCMON_HISTORY_SIZE)
        event_count++;

    spin_unlock_irqrestore(&event_lock, flags);
}
\end{lstlisting}

\subsection{Userland-Facing \file{/proc} Files}

The module publishes two read-only files. \file{stats} prints the cumulative
counters and the buffer fill level; \file{events} walks the ring buffer in
chronological order and prints a formatted table:

\begin{lstlisting}[language=C, caption={The events table walker (aulinux\_procmon.c, lines 132--158)}]
static int procmon_events_show(struct seq_file *m, void *v)
{
    unsigned long flags;
    int i, idx;
    static const char *type_names[] = { "FORK", "EXEC", "EXIT" };

    seq_printf(m, "Recent Process Events\n");
    seq_printf(m, "%-6s %-6s %-6s %-16s %s\n",
               "TYPE", "PID", "PPID", "COMM", "TIMESTAMP");

    spin_lock_irqsave(&event_lock, flags);
    for (i = 0; i < event_count; i++) {
        idx = (event_head - event_count + i + PROCMON_HISTORY_SIZE)
              % PROCMON_HISTORY_SIZE;
        seq_printf(m, "%-6s %-6d %-6d %-16s %lld\n",
                   type_names[event_history[idx].type],
                   event_history[idx].pid,
                   event_history[idx].ppid,
                   event_history[idx].comm,
                   ktime_to_ns(event_history[idx].timestamp));
    }
    spin_unlock_irqrestore(&event_lock, flags);
    return 0;
}
\end{lstlisting}

A userland \code{ps} could poll \file{/proc/aulinux/procmon/events} and join the
recent fork/exit stream against its own \file{/proc/[pid]} scan to flag
short-lived processes that a snapshot would miss.

\begin{notebox}[EXEC is declared but never recorded]
The \code{enum procmon\_event\_type} defines \code{PROCMON\_EXEC} and the table
header reserves an ``EXEC'' label and an \code{atomic\_long\_t total\_execs}
counter — but \emph{no kprobe targets exec}. There is no probe on
\code{do\_execve} / \code{bprm\_execve}, so \code{total\_execs} is permanently
zero and no EXEC rows ever appear. Exec monitoring is a stub. Likewise,
\code{monitoring\_enabled} is a hardcoded \code{static int = 1} with no
\file{/proc} or sysctl write handler, so the ``enable/disable'' capability
advertised in \file{stats} cannot actually be toggled at runtime yet.
\end{notebox}

\begin{warnbox}[A real bug: the procmon \file{/proc} directory]
\code{aulinux\_procmon\_init} calls
\code{proc\_mkdir("aulinux/procmon", NULL)} (line 209). This depends on
\file{/proc/aulinux} already existing — i.e. on \code{aulinux\_core} being
loaded \emph{first}. If procmon is loaded alone, the nested \code{proc\_mkdir}
fails and \code{procmon\_dir} is \code{NULL}; the code guards on this
(\code{if (procmon\_dir)}), so the \file{stats}/\file{events} files are simply
never created and the module loads ``successfully'' but exposes nothing. The
\file{Makefile} \code{load} target sidesteps this by ordering the loads
(\code{sysctl}, then \code{core}, then \code{procmon}), but there is no
\code{MODULE\_SOFTDEP} or symbol dependency enforcing it.
\end{warnbox}

\section{\code{aulinux\_sysctl}: Runtime Tunables}

\file{aulinux\_sysctl.c} registers the \file{/proc/sys/aulinux/} tree and is the
one module that includes the shared \file{aulinux.h}. It owns the three
configuration variables and exports the two accessor functions the header
promised.

\subsection{The \code{ctl\_table}}

Sysctl variables are described declaratively in a \code{struct ctl\_table}
array. Each entry binds a procfs name to a backing variable, a size, a mode, and
a handler that marshals between the text in \file{/proc/sys} and the in-kernel
integer. Two of the three entries use \code{proc\_dointvec\_minmax} with bound
variables to clamp input.

\begin{lstlisting}[language=C, caption={The sysctl table with min/max clamping (aulinux\_sysctl.c, lines 32--58)}]
static struct ctl_table aulinux_sysctl_table[] = {
    {
        .procname     = "debug_level",
        .data         = &aulinux_debug_level,
        .maxlen       = sizeof(int),
        .mode         = 0644,
        .proc_handler = proc_dointvec_minmax,
        .extra1       = &debug_level_min,   /* 0 */
        .extra2       = &debug_level_max,   /* 2 */
    },
    {
        .procname     = "feature_flags",
        .data         = &aulinux_feature_flags,
        .maxlen       = sizeof(int),
        .mode         = 0644,
        .proc_handler = proc_dointvec,      /* unbounded bitfield */
    },
    {
        .procname     = "max_processes",
        .data         = &aulinux_max_processes,
        .maxlen       = sizeof(int),
        .mode         = 0644,
        .proc_handler = proc_dointvec_minmax,
        .extra1       = &max_proc_min,      /* 128   */
        .extra2       = &max_proc_max,      /* 65536 */
    },
};
\end{lstlisting}

These files are mode \code{0644} — readable by all, writable by root. Writing
\code{echo 3 > /proc/sys/aulinux/debug\_level} is rejected by the
\code{minmax} handler (max is 2); \code{echo 2} succeeds and immediately changes
what \code{aulinux\_get\_debug\_level()} returns.

\subsection{Registration and the Exported Contract}

Registration uses the modern single-call \code{register\_sysctl()} API (the
older \code{register\_sysctl\_table} path is gone on current kernels). The module
then exports its two accessors so that other modules — and the
\code{aulinux\_pr\_debug} macro in the shared header — can query configuration
at runtime.

\begin{lstlisting}[language=C, caption={Exported symbols and sysctl registration (aulinux\_sysctl.c, lines 65--91)}]
int aulinux_get_debug_level(void)
{
    return aulinux_debug_level;
}
EXPORT_SYMBOL(aulinux_get_debug_level);

int aulinux_feature_enabled(int flag)
{
    return (aulinux_feature_flags & flag) != 0;
}
EXPORT_SYMBOL(aulinux_feature_enabled);

static int __init aulinux_sysctl_init(void)
{
    aulinux_sysctl_header = register_sysctl("aulinux", aulinux_sysctl_table);
    if (!aulinux_sysctl_header) {
        pr_err("AULinux: Failed to register sysctl table\n");
        return -ENOMEM;
    }
    return 0;
}
\end{lstlisting}

\begin{notebox}[The exports are real, the consumers are mostly latent]
\code{EXPORT\_SYMBOL} genuinely publishes both functions into the kernel symbol
table, so this module \emph{must} be loaded before any module that references
them (hence its first position in the \file{Makefile} \code{load} order).
However, today only \file{aulinux.h}'s \code{aulinux\_pr\_debug} macro would call
\code{aulinux\_get\_debug\_level()}, and that macro is not yet used in
\file{aulinux\_core.c} or \file{aulinux\_procmon.c}. \code{aulinux\_feature\_enabled()}
and the \code{max\_processes} tunable have no in-kernel readers at all — they are
plumbed and ready, but nothing consumes them yet.
\end{notebox}

\section{The Three Modules Side by Side}

\begin{table}[h]
\centering
\begin{tabular}{@{}llll@{}}
\toprule
\textbf{Module} & \textbf{/proc or /sys entry} & \textbf{Purpose} & \textbf{Status} \\
\midrule
\code{aulinux\_core}    & \file{/proc/aulinux/version}      & Version telemetry            & Complete \\
\code{aulinux\_procmon} & \file{/proc/aulinux/procmon/stats}  & Lifecycle counters         & Partial (exec stubbed) \\
                        & \file{/proc/aulinux/procmon/events} & Event ring buffer          & Fork probe broken on 5.10+ \\
\code{aulinux\_sysctl}  & \file{/proc/sys/aulinux/*}        & Runtime tunables + exports   & Complete; consumers latent \\
\bottomrule
\end{tabular}
\caption{Real entry points and honest status of the AULinux kernel modules.}
\end{table}

\section{Building and Loading}

The modules are built with the standard kbuild out-of-tree pattern. Both
\file{src/Kbuild} and the top-level \file{Makefile} list the three objects as
\code{obj-m} targets, which tells kbuild to compile them as loadable modules
(\code{.ko}) rather than linking them into the kernel.

\begin{lstlisting}[language=C, caption={Module targets in src/Kbuild}]
obj-m += aulinux_core.o
obj-m += aulinux_sysctl.o
obj-m += aulinux_procmon.o

ccflags-y := -I$(src)
ccflags-y += -DAULINUX_MODULE
\end{lstlisting}

The \file{Makefile} drives the kernel's own build system by recursion. The key
invocation hands control to the kernel tree at \code{KDIR} and points it back at
the AULinux sources via \code{M=}:

\begin{lstlisting}[language=bash, caption={The out-of-tree build recursion (Makefile)}]
KDIR ?= /lib/modules/$(shell uname -r)/build
SRC_DIR := $(shell pwd)/src

modules:
	$(MAKE) -C $(KDIR) M=$(SRC_DIR) modules
	@cp $(SRC_DIR)/*.ko $(BUILD_DIR)/    # -> ../build/kernel/
\end{lstlisting}

This is why the modules need the matching \textbf{kernel headers / build
directory} installed: kbuild compiles each \code{.c} against the exact
configuration of the target kernel and stamps the result with a vermagic string.
A \code{.ko} built for one kernel will refuse to load on a mismatched one.

\subsection{Load Order Matters}

The \file{Makefile} provides \code{load} and \code{unload} convenience targets
that encode the dependency ordering discovered above:

\begin{lstlisting}[language=bash, caption={Ordered insmod / rmmod (Makefile)}]
load:
	insmod $(BUILD_DIR)/aulinux_sysctl.ko    # exports symbols first
	insmod $(BUILD_DIR)/aulinux_core.ko      # creates /proc/aulinux
	insmod $(BUILD_DIR)/aulinux_procmon.ko   # nests under it

unload:
	-rmmod aulinux_procmon
	-rmmod aulinux_core
	-rmmod aulinux_sysctl
\end{lstlisting}

\begin{warnbox}[Why the order is not optional]
\code{aulinux\_sysctl} must load first because it \code{EXPORT\_SYMBOL}s the
accessors other modules link against. \code{aulinux\_core} must precede
\code{aulinux\_procmon} because the latter nests its directory under
\file{/proc/aulinux}. Because nothing in the code declares these dependencies
(no \code{MODULE\_SOFTDEP}, no referenced exported symbol from procmon), loading
out of order fails silently rather than loudly. In a robust build this ordering
should be promoted into \code{depmod} metadata or explicit soft dependencies.
\end{warnbox}

\section{Design Rationale and Limitations}

The AULinux kernel layer is a clean, readable demonstration of the three pillars
of kernel-userspace interaction — \code{procfs} via \code{seq\_file},
\code{sysctl} via \code{ctl\_table}, and dynamic instrumentation via kprobes —
but it is honestly an early-stage scaffold rather than a finished subsystem.

\begin{itemize}
  \item \textbf{\code{aulinux\_core}} is complete and correct: it does exactly
        one thing (publish version telemetry) and does it idiomatically with
        proper failure unwinding.
  \item \textbf{\code{aulinux\_sysctl}} is complete as a control surface — the
        tunables clamp, persist, and are exported — but its readers are still
        latent, so today it configures behaviour nothing yet acts on.
  \item \textbf{\code{aulinux\_procmon}} is the furthest from done. Its fork
        probe targets a symbol (\code{do\_fork}) removed from modern kernels,
        exec monitoring is declared but unimplemented, and
        \code{monitoring\_enabled} cannot be toggled. The ring-buffer plumbing
        and the \file{/proc} presentation, however, are real and sound.
\end{itemize}

\begin{infobox}[The deeper lesson]
Kernel code runs with no memory protection and no safety net: a stray pointer in
a probe handler panics the whole machine. AULinux's choice to keep these modules
small, single-purpose, and loadable — rather than patching the monolith — is
precisely the discipline that keeps that danger contained. The honest gaps
catalogued above are not failures of the design; they are the to-do list a
serious kernel engineer would tackle next, with the architecture already laid
out to receive them.
\end{infobox}
